\documentclass[11pt,a4paper]{article}
\usepackage[margin=1in]{geometry}
\usepackage{amsmath,amssymb,amsfonts}
\usepackage{mathtools}
\usepackage{lmodern}
\usepackage[T1]{fontenc}
\usepackage[utf8]{inputenc}
\usepackage{textcomp}
\usepackage{microtype}
\usepackage{parskip}
\usepackage{enumitem}
\usepackage{array}
\usepackage{booktabs}
\usepackage{hyperref}
\hypersetup{hidelinks,
  pdftitle={The V5 Finite-Memory Kernel as a Substrate-Level Resolution of the Trans-Planckian Problem},
  pdfauthor={Allen Proxmire}}
\setlength{\parindent}{0pt}
\setlength{\parskip}{6pt plus 2pt minus 1pt}

\title{\vspace{-1cm}\Large\bfseries The V5 Finite-Memory Kernel as a Substrate-Level Resolution of the Trans-Planckian Problem}
\author{Allen Proxmire \\ \textit{Independent Researcher}}
\date{May 2026}

\begin{document}
\maketitle

\section*{Abstract}

The standard semiclassical derivation of Hawking radiation traces observed-at-infinity modes back to arbitrarily-blueshifted near-horizon frequencies. The derivation thus implicitly requires standard QFT to apply at and beyond the Planck scale, where quantum-gravity effects should structurally modify field theory. This is the \emph{trans-Planckian problem}, recognized since Jacobson [1] (1991) and analyzed extensively in subsequent literature [2--6]. Multiple proposals address it phenomenologically: modified dispersion relations [1, 2], lattice cutoffs [3], naturalness arguments [4], the Trans-Planckian Censorship Conjecture (TCC) [5], and others. We show that the Event Density (ED) framework's V5 finite-memory cross-chain kernel, with characteristic memory time $\tau_{V5} = \ell_P/c$ at the gravitational scale (the natural substrate-built timescale via T19 + dimensional analysis), produces a substrate-level resolution that does not require any of the standard phenomenological additions. The substrate cannot mediate cross-chain correlations at frequencies $\omega \gg c/\ell_P$ because V5's coherence breaks down beyond the cutoff. Combined with the Diffusion Coarse-Graining Theorem's hydrodynamic-window closure at $L_\mathrm{flow} \sim \ell_P$, the substrate naturally cuts off arbitrarily-blueshifted modes near the horizon. The standard trans-Planckian divergence is a continuum artifact that vanishes at the substrate level. The V5-modulated Hawking spectrum is $N_{ED}(\omega) = N_H(\omega) / (1+(\omega\tau_{V5})^{2})$, finite at all frequencies and reducing to standard Hawking at observable scales. Falsifiable predictions in current-generation analog Hawking experiments and primordial-BH late-stage evaporation distinguish the framework from modified-dispersion and lattice-cutoff approaches.

\textbf{Keywords:} trans-Planckian problem, Hawking radiation, Event Density, substrate ontology, V5 finite-memory kernel, UV regulation, modified dispersion, Diffusion Coarse-Graining Theorem.

\section{Introduction}

In the standard semiclassical derivation of Hawking radiation [7, 8], a Schwarzschild black hole of mass $M$ emits thermal radiation at temperature $T_H = \kappa/(2\pi)$, where $\kappa$ is the surface gravity of the horizon. The derivation runs through QFT in curved spacetime: vacuum modes propagated from past null infinity, through the BH-collapse spacetime, to future null infinity, where they appear as a thermal mixture via the Bogoliubov transformation between past and future mode bases.

The derivation has empirical content (analog Hawking experiments [9--10] confirm the spectral form) but a structural concern. Modes observed at moderate frequencies $\omega \sim T_H$ at substrate-asymptotic infinity were arbitrarily-blueshifted near the horizon at proper frequencies
\[
\omega_\mathrm{proper}(r) = \omega \cdot (1 - 2M/r)^{-1/2}
\]
For $r \to r_h^+$ (horizon-approaching), $\omega_\mathrm{proper} \to \infty$. The standard derivation thus implicitly requires standard QFT to apply at proper frequencies arbitrarily far above the Planck scale $\omega_P = c/\ell_P \approx 1.85 \times 10^{43}$ Hz, where quantum-gravity effects should structurally modify field theory.

This is the \emph{trans-Planckian problem} of Hawking radiation [1]. It has been addressed in the literature through several proposals:

\begin{itemize}[noitemsep]
\item \textbf{Modified dispersion relations} [1, 2]: postulate UV-modified dispersion $\omega^{2} = c^{2} k^{2} (1 + \mathcal{O}(\ell_P k))$ at the Planck scale.
\item \textbf{Lattice cutoffs} [3]: discretize spacetime at the Planck scale, eliminating modes with $k \gg 1/\ell_P$.
\item \textbf{Naturalness / robustness arguments} [4]: argue that physics-domain robustness should protect Hawking radiation from trans-Planckian modifications.
\item \textbf{Trans-Planckian Censorship Conjecture (TCC)} [5]: conjectured cosmological censorship of trans-Planckian modes by sufficient inflation.
\item \textbf{'t Hooft S-matrix arguments} [6]: structural arguments from BH information theory.
\end{itemize}

Each proposal is structurally meaningful but each is \emph{phenomenological} in character: each posits the cutoff or modification rather than deriving one from a substrate-level ontology. The trans-Planckian problem has therefore remained a structural concern of standard semiclassical Hawking, without a derivation-level resolution.

This paper presents the Event Density (ED) framework's substrate-level resolution. The framework is built on a substrate ontology of discrete micro-events, finite participation bandwidth, irreversible commitment events, and finite-width substrate kernels [11, 12]. The framework's V5 finite-memory cross-chain kernel --- a primitive substrate object with characteristic memory time $\tau_{V5}$ --- naturally regulates the high-frequency behavior of substrate-mediated correlations. At the gravitational scale, $\tau_{V5} = \ell_P/c$ (the Planck time, identified via T19 + substrate dimensional analysis [13]). This is the natural substrate-built timescale at which V5 coherence breaks down.

The framework's claim: the substrate cannot mediate cross-chain correlations at frequencies $\omega \gg c/\ell_P$. Arbitrarily-blueshifted near-horizon modes do not exist at the substrate level. The trans-Planckian divergence of standard semiclassical Hawking is a \emph{continuum artifact} of the QFT-in-curved-spacetime calculation; the substrate-level account regulates it naturally without postulating modified dispersion, lattice cutoffs, or any other phenomenological addition.

Combined with the Diffusion Coarse-Graining Theorem's (DCGT) hydrodynamic-window closure at $L_\mathrm{flow} \sim \ell_P$, the framework provides a substrate-level UV regulation that:
\begin{enumerate}[noitemsep]
\item Reproduces standard Hawking exactly at observable scales.
\item Cuts off the high-frequency tail at the Planck frequency.
\item Resolves the trans-Planckian problem without new physics or postulated regulators.
\item Produces falsifiable first-subleading-order corrections distinguishing the framework from standard semiclassical Hawking and from the major phenomenological resolutions.
\end{enumerate}

The structure of the paper: \S2 states the standard trans-Planckian problem precisely. \S3 presents the V5 kernel at the substrate level. \S4 derives the V5-modulated Hawking spectrum. \S5 articulates DCGT's role in the resolution. \S6 states the substrate-level resolution explicitly. \S7 contrasts with the major standard-physics proposals. \S8 identifies falsifiable predictions. \S9 discusses the implications.

\section{The Standard Trans-Planckian Problem}

\subsection{The blueshift argument}

In Schwarzschild geometry, an outgoing mode observed at substrate-asymptotic infinity with frequency $\omega$ traces back to a mode near the horizon at proper frequency
\[
\omega_\mathrm{proper}(r) = \omega \cdot (1 - 2M/r)^{-1/2}.
\]
For an observer at radius $r = 2M(1+\epsilon)$ with small $\epsilon$, the proper frequency is approximately $\omega/\sqrt{2\epsilon}$. As $\epsilon \to 0$, $\omega_\mathrm{proper} \to \infty$. Modes traced back to the horizon at any positive distance away are blueshifted to arbitrarily high proper frequencies.

For a stellar-mass black hole observed mode at $\omega \sim T_H$ (peak of the Hawking spectrum), the proper frequency exceeds the Planck scale at distances
\[
\epsilon \lesssim (T_H/\omega_P)^{2} \cdot 2 \sim 10^{-86}
\]
--- a substantial fraction of the horizon's neighborhood, in the relevant blueshift sense.

\subsection{The structural concern}

Standard QFT in curved spacetime is a continuum-level theory whose validity at and beyond the Planck scale is structurally questionable. Quantum-gravity effects should modify field theory at the Planck scale, where:
\begin{itemize}[noitemsep]
\item Spacetime fluctuations become order-unity.
\item Quantum-gravity corrections to dispersion relations should appear.
\item The continuum description of spacetime should break down.
\end{itemize}

The trans-Planckian problem articulates this concern: the standard semiclassical Hawking calculation requires QFT to be valid at arbitrarily-high proper frequencies near the horizon, including frequencies far above where QFT's structural status is secure.

\subsection{The standard responses}

Multiple proposals have been advanced:

\textbf{(a) Modified dispersion relations [1, 2].} Postulate that the standard relativistic dispersion $\omega^{2} = c^{2} k^{2}$ is modified at the Planck scale to $\omega^{2} = c^{2} k^{2} + f(k\ell_P)$ for some function $f$ that effectively cuts off mode propagation at $k \sim 1/\ell_P$. This requires postulating the specific form of $f$.

\textbf{(b) Lattice cutoffs [3].} Discretize spacetime at the Planck scale. This requires postulating the discrete spacetime structure and its specific properties.

\textbf{(c) Naturalness / robustness arguments [4].} Argue that the empirical robustness of Hawking radiation should protect it from trans-Planckian sensitivity. This is an argument-from-empirical-robustness rather than a derivation.

\textbf{(d) Trans-Planckian Censorship Conjecture (TCC) [5].} Conjecture that cosmological inflation must terminate before trans-Planckian modes have stretched to observable scales. This is a conjecture about cosmological dynamics; it does not directly resolve the BH trans-Planckian problem.

\textbf{(e) 't Hooft S-matrix arguments [6].} Structural arguments from BH information theory and the holographic principle, suggesting that the BH S-matrix should be unitary and structurally regulated.

Each proposal contributes structural insight; each remains phenomenological in character. The trans-Planckian problem awaits a derivation-level resolution.

\section{The V5 Finite-Memory Kernel}

\subsection{Substrate definition}

The V5 cross-chain memory kernel is a primitive substrate object in the Event Density framework [11, 12]. Its primitive structural form is determined by:

\begin{itemize}[noitemsep]
\item \textbf{Forward-cone support.} V5 mediates cross-chain correlations causally forward in substrate-time (inherited from T18 [14]).
\item \textbf{Finite memory.} V5 has a characteristic correlation time $\tau_{V5}$ --- a substrate-determined timescale that varies by physical context.
\item \textbf{Substrate locality.} V5 mediates correlations only along substrate-locality-permitted pathways.
\end{itemize}

The minimal substrate-consistent kernel form is:
\[
V_5(t) = \mathcal{V}_0 \cdot \theta(t) \cdot e^{-t/\tau_{V5}} \cdot \psi(t/\tau_{V5})
\]
where $\theta(t)$ is the Heaviside step function (forward-cone-only), $e^{-t/\tau_{V5}}$ is the leading-order decay envelope, $\psi(t/\tau_{V5})$ is a substrate-determined dimensionless shape function approaching unity at small argument, and $\mathcal{V}_0$ is the substrate-coupling normalization. The exact functional form of $\psi$ is INHERITED from substrate-microscopic details and is not fixed at the framework's structural-foundations level.

\subsection{Frequency-domain structure}

The Fourier transform of the simplest substrate-consistent V5 form ($\psi = 1$):
\[
\tilde V_5(\omega) = \frac{\mathcal{V}_0 \tau_{V5}}{1 - i\omega\tau_{V5}}
\]
The magnitude squared, governing V5's coherent transfer of cross-chain correlations:
\[
|\tilde V_5(\omega)|^{2} = \frac{(\mathcal{V}_0 \tau_{V5})^{2}}{1 + (\omega\tau_{V5})^{2}}
\]
For $\omega \ll \omega_c \equiv 1/\tau_{V5}$, $|\tilde V_5(\omega)|^{2}$ is approximately constant --- V5 is effectively local-in-time and behaves as a delta-function memory.

For $\omega \gtrsim \omega_c$, $|\tilde V_5(\omega)|^{2} \to (1/\omega^{2})$ --- V5 coherence breaks down. The substrate cannot mediate cross-chain correlations at frequencies above the cutoff scale.

\subsection{The cutoff scale at the gravitational scale}

The cutoff $\tau_{V5}$ is a substrate-determined parameter. At the gravitational scale, the only substrate-built timescale available from substrate primitives is the Planck time:
\[
\tau_{V5} = \ell_P/c \approx 5.4 \times 10^{-44}\, \mathrm{s}
\]
This identification is FORCED via T19 (Newton-recovery $\ell_P$) [13] plus substrate dimensional analysis. The Planck time is the natural substrate timescale at which V5 coherence breaks down at gravitational scales.

The corresponding cutoff frequency:
\[
\omega_c = c/\ell_P \approx 1.85 \times 10^{43}\, \mathrm{Hz}
\]
--- the Planck frequency. The V5 kernel does not coherently mediate cross-chain correlations at proper frequencies approaching the Planck scale.

\subsection{Why $\tau_{V5}$ is not phenomenological}

The V5 cutoff scale is \emph{not} a phenomenological cutoff inserted ad hoc to regulate the trans-Planckian problem. The V5 kernel is a substrate primitive in its own right --- it appears across multiple framework sectors:
\begin{itemize}[noitemsep]
\item Soft-matter Maxwell viscoelasticity, where $\tau_{V5}$ identifies with molecular relaxation times [15].
\item Hawking radiation high-frequency spectral cutoff (this paper).
\item Black-hole information transfer bandwidth modulation [16, 17].
\item Substrate-Unruh radiation cutoff [18].
\end{itemize}

The kernel exists \emph{before} and \emph{independently of} the trans-Planckian problem. The cutoff scale at the gravitational context is FORCED by substrate dimensional analysis. The trans-Planckian regulation is a \emph{consequence} of V5's primitive structure, not a feature added to address the problem.

\section{The V5-Modulated Hawking Spectrum}

\subsection{The leading-order substrate Hawking calculation}

The framework's leading-order substrate-level Hawking calculation [19, 20] reproduces the standard semiclassical result via the substrate-Unruh argument applied at a saturated decoupling surface (substrate-level analog of a BH horizon). V5 cross-chain correlations across the surface, satisfying imaginary-time periodicity at $\beta = 2\pi/\kappa_{ED}$ (forced by substrate-vacuum self-consistency at the surface), produce a Planck spectrum at temperature $T = \kappa_{ED}/(2\pi)$. DCGT identification at leading order gives $T_H = \kappa/(2\pi)$ --- the standard Hawking temperature.

The leading-order spectrum:
\[
N_H(\omega) = \frac{1}{e^{\omega/T_H} - 1}
\]
This is the standard result, exactly reproduced.

\subsection{First-subleading-order V5 modulation}

At first subleading order, V5's finite-memory structure modulates the spectrum. The substrate-level near-horizon correlation function involves V5 cross-chain mediation at frequency $\omega$. The leading-order calculation treats V5 as effectively local-in-time (valid for $\omega \ll \omega_c$). The first-subleading correction includes the explicit V5 frequency-domain modulation:
\[
N_{ED}(\omega) = N_H(\omega) \cdot \frac{|\tilde V_5(\omega)|^{2}}{|\tilde V_5(0)|^{2}} = \frac{N_H(\omega)}{1 + (\omega\tau_{V5})^{2}}
\]
This is the V5-modulated Hawking spectrum. For modes with $\omega \ll \omega_c = c/\ell_P$:
\[
N_{ED}(\omega) \approx N_H(\omega) \cdot [1 - (\omega\tau_{V5})^{2} + \mathcal{O}((\omega\tau_{V5})^{4})]
\]
The modulation is approximately unity, and the spectrum reduces to the standard Planck distribution.

For modes with $\omega \gg \omega_c$:
\[
N_{ED}(\omega) \approx N_H(\omega) \cdot (\omega_c/\omega)^{2}
\]
The modulation suppresses the high-frequency tail. The substrate-level spectrum is finite and well-defined at all frequencies.

\subsection{The peak of the spectrum}

The Hawking spectrum's peak is at $\omega \sim T_H \cdot k_B/\hbar$. For a Schwarzschild stellar-mass BH:
\[
T_H \approx 6 \times 10^{-8}\, \mathrm{K}, \qquad \omega_\mathrm{peak} \sim 10^{4}\, \mathrm{Hz}
\]
The corresponding $\omega_\mathrm{peak}\tau_{V5} \sim 10^{4} \cdot 10^{-43} \sim 10^{-39}$ --- invisible at the spectrum peak.

For a primordial BH in its final stages of evaporation ($M \to M_P$):
\[
T_H \to T_P \sim 10^{32}\, \mathrm{K}, \qquad \omega_\mathrm{peak} \sim \omega_c
\]
The V5 cutoff becomes order-unity at the peak. The spectrum is dramatically modified from standard Hawking at extreme scales.

\subsection{Integrated emission rate}

The framework's integrated Page emission rate is also modulated by V5 [21]:
\[
\left.\frac{dM}{dt}\right|_{ED} = -\frac{\alpha_\mathrm{Page}}{M^{2}} \cdot \left[1 - K\left(\frac{\ell_P}{M}\right)^{2} + \mathcal{O}\left(\frac{\ell_P}{M}\right)^{4}\right]
\]
where $K = -c_{V5}^\mathrm{int} - c_g^\mathrm{int}\log g$ is an order-unity coefficient inherited from V5 and motif-alphabet substrate details. The correction is invisible at observable BH scales but becomes order-unity at $M \sim M_P$.

\section{DCGT and the Closing Hydrodynamic Window}

The V5 cutoff is one structural component of the framework's substrate-level resolution. The Diffusion Coarse-Graining Theorem (DCGT) [15] supplies the second.

\subsection{The DCGT scale-separation requirement}

DCGT establishes the substrate-to-continuum bridge through hydrodynamic-window scale separation:
\[
\ell_P \ll R_\mathrm{cg} \ll L_\mathrm{flow}
\]
The substrate-level lower bound is $\ell_P$ (the substrate's irreducible length scale). The upper bound $L_\mathrm{flow}$ is the characteristic length scale of the continuum theory under derivation. For Hawking radiation, $L_\mathrm{flow}$ is set by the BH's size: $L_\mathrm{flow} = r_s = 2M$.

The scale separation requires $M \gg \ell_P$ substantially. For ordinary BHs ($M \gg M_P$), the window is enormous (ratios of $10^{38}$ to $10^{60}$ are typical). For Planck-mass-approaching BHs ($M \sim M_P$), the window closes.

\subsection{The window closure and its meaning}

When $M$ approaches $\ell_P$, the upper and lower bounds of the hydrodynamic window coincide. There is no scale $R_\mathrm{cg}$ satisfying $\ell_P \ll R_\mathrm{cg} \ll 2M$. The DCGT substrate-to-continuum bridge fails.

What does this mean structurally? Hawking radiation as a continuum-level phenomenon depends on the DCGT bridge. The bridge's leading-order consequences (scalar diffusion, directional viscosity, V1$\to$R1 substrate-cutoff, V5$\to$Maxwell viscoelasticity, T17 minimal coupling [15]) all require the hydrodynamic-window separation. When the window closes, the substrate-level dynamics no longer produce continuum-level emission as a coherent process.

For the trans-Planckian problem specifically: blueshifted near-horizon modes at proper frequencies $\omega_\mathrm{proper} \gtrsim 1/\ell_P$ have proper wavelengths $\lambda_\mathrm{proper} \lesssim \ell_P$. These modes lie \emph{below} the DCGT hydrodynamic window's lower bound. They are not continuum-level features at all.

\subsection{The combined substrate-level cutoff}

The two structural components combine to provide the substrate-level cutoff:

\begin{enumerate}[noitemsep]
\item \textbf{V5 finite-memory cutoff (perturbative).} The V5 kernel's coherent response saturates at $\omega \sim 1/\tau_{V5} = c/\ell_P$. Substrate cross-chain correlations at frequencies above the cutoff are suppressed by $1/\omega^{2}$. This is the leading-order substrate-cutoff structure.
\item \textbf{DCGT hydrodynamic-window closure (non-perturbative).} When proper wavelengths approach $\ell_P$, the substrate-to-continuum bridge fails entirely. The continuum description does not apply. This is the structural-foundational substrate-cutoff statement.
\end{enumerate}

The two components are complementary. V5 provides a smooth perturbative cutoff at frequencies below the DCGT failure scale. DCGT provides the structural-foundational statement that the continuum description does not extend into the trans-Planckian regime.

\section{The Substrate-Level Resolution}

\subsection{Statement of the resolution}

\textbf{The substrate-level trans-Planckian resolution.} Substrate cross-chain correlations at frequencies $\omega \gg c/\ell_P$ are not coherently mediated. The substrate's V5 finite-memory kernel does not support coherent correlations beyond $\tau_{V5}^{-1} = c/\ell_P$. Combined with DCGT's hydrodynamic-window closure when proper wavelengths approach $\ell_P$, the substrate provides a natural UV cutoff at the Planck frequency. Modes at proper frequencies above the Planck scale do not exist as continuum-level objects. The standard trans-Planckian divergence of semiclassical Hawking is a \emph{continuum artifact} of extending QFT beyond its substrate-level support; the substrate-level account regulates the divergence naturally without postulating modified dispersion, lattice cutoffs, or any phenomenological addition.

\subsection{Why this is not a phenomenological cutoff}

The framework's resolution is structurally distinct from postulating an ad-hoc cutoff:

\begin{itemize}[noitemsep]
\item The V5 kernel is a substrate primitive that exists \emph{independently} of the trans-Planckian problem. The same kernel produces Maxwell viscoelastic memory in soft matter and entanglement-bandwidth modulation in BH information transfer [16].
\item The cutoff scale $\tau_{V5} = \ell_P/c$ at the gravitational scale is FORCED by T19 + substrate dimensional analysis. It is the unique substrate-built timescale at the gravitational context.
\item DCGT's hydrodynamic-window closure is a structural feature of the substrate-to-continuum bridge, not a regulator inserted to address the trans-Planckian problem.
\end{itemize}

The trans-Planckian regulation is therefore a \emph{consequence} of substrate ontology, not a feature added to address a structural concern. This is the methodological distinction between phenomenological regularization and substrate-level resolution.

\subsection{What this delivers}

The framework's resolution:
\begin{enumerate}[noitemsep]
\item Reproduces standard Hawking exactly at observable scales.
\item Produces a finite, well-defined Hawking spectrum at all frequencies.
\item Cuts off the high-frequency tail at the Planck frequency via V5.
\item Establishes the structural-foundational cutoff via DCGT closure when wavelengths approach $\ell_P$.
\item Predicts substrate-cutoff corrections at first subleading order $(\ell_P/M)^{2}$.
\item Combined with H-8's higher-order resummation [21], produces a stable Planck-mass remnant at the late-time evaporation endpoint.
\end{enumerate}

The resolution is therefore not just a regularization --- it is a substrate-level account of where, why, and how the standard semiclassical framework's structural concerns are resolved.

\section{Comparison with Standard Proposals}

\subsection{Modified dispersion relations (Unruh, Jacobson, BMPS)}

Modified-dispersion approaches [1, 2] postulate UV-modified dispersion relations such as:
\[
\omega^{2} = c^{2} k^{2} + f(k\ell_P)
\]
with $f(k\ell_P)$ a function that effectively cuts off mode propagation at $k \sim 1/\ell_P$. Specific forms include $f(k\ell_P) = -(c\ell_P k)^{4}/c^{2}$ (Unruh-class), and various lattice-derived forms.

\textbf{Comparison with ED:}
\begin{itemize}[noitemsep]
\item \emph{Modified dispersion}: postulates the dispersion modification as a phenomenological feature.
\item \emph{ED}: derives the high-frequency cutoff from V5's substrate primitive structure. The dispersion relation itself is unchanged at the leading order; the modification appears as a multiplicative modulation of the spectrum, not as a modification of the underlying field-theoretic mode structure.
\end{itemize}

\textbf{Empirical distinguishability:} Modified-dispersion approaches typically predict specific UV modification forms that depend on the chosen $f(k\ell_P)$. ED's V5 modulation $1/(1+(\omega\tau_{V5})^{2})$ is form-FORCED by V5's substrate primitive. Precision analog Hawking experiments could distinguish between different cutoff functional forms.

\subsection{Lattice cutoffs}

Lattice-cutoff approaches [3] discretize spacetime at the Planck scale, eliminating modes with wavelengths shorter than the lattice spacing.

\textbf{Comparison with ED:}
\begin{itemize}[noitemsep]
\item \emph{Lattice cutoff}: discretizes spacetime structure at the Planck scale by postulate.
\item \emph{ED}: substrate ontology is \emph{already} discrete (P01: discrete micro-events). The substrate ontology underlying ED is not a regularization of continuum spacetime but the framework's primitive ontology. The continuum description emerges as a coarse-grained leading-order consequence (DCGT [15]).
\end{itemize}

\textbf{Structural distinction:} ED's discreteness is ontologically primary; lattice cutoff approaches treat discreteness as a regulator added to a continuum-level theory. In ED, the continuum is the emergent description; in lattice approaches, the continuum is the fundamental description being regulated.

\subsection{Naturalness and robustness arguments}

Naturalness arguments [4] argue that Hawking radiation should be \emph{robust} against trans-Planckian sensitivities --- that the empirically validated content (the thermal spectrum at $T_H = \kappa/(2\pi)$) is protected from quantum-gravity modifications by general arguments about effective field theory.

\textbf{Comparison with ED:}
\begin{itemize}[noitemsep]
\item \emph{Robustness arguments}: argue from empirical-protection that trans-Planckian content shouldn't matter, without supplying a mechanism.
\item \emph{ED}: provides the explicit mechanism (V5 substrate-cutoff + DCGT closure) by which trans-Planckian modes are absent at the substrate level.
\end{itemize}

\textbf{Methodological distinction:} Robustness arguments are arguments-from-empirical-protection; ED supplies the underlying mechanism that makes the empirical robustness structurally inevitable.

\subsection{Trans-Planckian Censorship Conjecture (TCC)}

TCC [5] conjectures that cosmological inflation must terminate before trans-Planckian modes have stretched to observable scales. This addresses cosmological trans-Planckian sensitivity but not directly the BH trans-Planckian problem.

\textbf{Comparison with ED:}
\begin{itemize}[noitemsep]
\item \emph{TCC}: cosmological constraint on inflation duration.
\item \emph{ED}: substrate-level UV cutoff applicable in any context (BH, cosmological horizon, accelerated observer [18]).
\end{itemize}

\textbf{Scope distinction:} TCC and ED's V5 cutoff address different aspects of the trans-Planckian problem. TCC concerns cosmological inflation; ED's V5 cutoff applies to all substrate-level modes. The two could in principle coexist; ED's V5 mechanism does not require TCC.

\subsection{Firewalls and information-paradox-class proposals}

Firewall [22] and related proposals (ER=EPR [23], soft hair [24], island formula [25]) address BH information at the horizon scale but do not directly resolve the trans-Planckian problem. They address what happens \emph{at} or \emph{near} the horizon, not the structural question of whether QFT applies at trans-Planckian frequencies.

\textbf{Comparison with ED:} ED addresses a structurally different question. Firewalls / ER=EPR / soft hair / island formula address the BH information question (handled separately in [16]). The trans-Planckian problem is addressed by V5 + DCGT, independently of the BH information mechanism.

\subsection{Comparison summary}

\begin{center}
\begin{tabular}{p{0.18\textwidth}p{0.22\textwidth}p{0.18\textwidth}p{0.28\textwidth}}
\toprule
\textbf{Approach} & \textbf{Mechanism} & \textbf{Status} & \textbf{Empirical distinguishability} \\
\midrule
Modified dispersion & Postulated UV dispersion modification & Phenomenological & Specific cutoff form depends on chosen $f$ \\
Lattice cutoffs & Postulated spacetime discreteness & Phenomenological regulator & Discrete-spectrum signatures at small scales \\
Naturalness arguments & Empirical-protection argument & Argumentative & Cannot directly distinguish \\
TCC & Cosmological constraint on inflation & Conjectural & Cosmological inflation observables \\
ED V5 + DCGT & Substrate primitive + structural closure & Substrate-derived & Form-FORCED $1/(1+(\omega\tau)^{2})$ cutoff \\
\bottomrule
\end{tabular}
\end{center}

\section{Falsifiable Predictions}

The framework's substrate-level resolution distinguishes itself from standard-physics approaches through specific falsifiable predictions.

\subsection{Analog Hawking cutoff form}

In analog Hawking experiments [9, 10], the Hawking-analog spectrum should exhibit a high-frequency cutoff at the analog system's microscopic correlation timescale. The framework predicts:

\begin{itemize}[noitemsep]
\item The cutoff form is $1/(1+(\omega\tau_\mathrm{analog})^{2})$.
\item The analog timescale $\tau_\mathrm{analog}$ identifies with the analog system's characteristic correlation time:
\begin{itemize}[noitemsep]
\item BEC: $\tau_\mathrm{analog} \sim \xi/c_s$ (healing length over speed of sound).
\item Acoustic: $\tau_\mathrm{analog} \sim \lambda_\mathrm{phonon}/c_s$.
\item Photonic: $\tau_\mathrm{analog} \sim$ optical-cycle period.
\end{itemize}
\end{itemize}

Existing analog Hawking experiments confirm spectral form at moderate frequencies; precision tests at the cutoff scale are technically feasible with current-generation experiments. Modified-dispersion approaches predict specific cutoff functional forms depending on the chosen $f(k\ell_P)$; ED's prediction is the specific V5-derived $1/(1+(\omega\tau)^{2})$ form.

\subsection{Primordial BH late-stage spectra}

For primordial black holes in their final stages of evaporation, the framework predicts:
\begin{itemize}[noitemsep]
\item Spectral cutoff at $\omega \sim \omega_c = c/\ell_P$ (Planck frequency).
\item Modified evaporation profile slowing in the final $\sim 0.1$ s.
\item Possible Planck-mass remnant after evaporation halts (per H-8 [21]).
\end{itemize}

Standard semiclassical Hawking predicts complete evaporation in finite time with no upper-frequency cutoff. PBH detection with sufficient temporal and spectral resolution would test the framework's substrate-cutoff prediction.

\subsection{High-energy photon dispersion in vacuum}

If the V5 substrate-cutoff at $\omega_c = c/\ell_P$ extends to vacuum-substrate propagation (not only to BH horizon contexts), high-energy astrophysical photons traveling astronomical distances could exhibit dispersion at order $(\omega/\omega_c)^{2}$:
\[
\Delta v / c \sim (\omega/\omega_c)^{2}
\]
Public LIGO/Virgo + Fermi-LAT data on high-energy gamma-ray bursts and gravitational-wave-electromagnetic counterparts permit precision time-of-flight measurements. This connects ED's V5 substrate-cutoff to the active observational program on Lorentz-invariance-violation searches at high energies [26].

\subsection{Cross-platform consistency}

The V5 kernel does multiple substrate-level jobs across the framework's other domains:
\begin{itemize}[noitemsep]
\item Soft-matter Maxwell viscoelasticity at molecular relaxation timescales.
\item Substrate-Unruh radiation cutoff at accelerated-observer Rindler-like horizons.
\item Black-hole information transfer bandwidth modulation.
\end{itemize}

Cross-platform measurement of V5-mediated phenomena across these distinct domains provides a structural consistency test. If the V5 kernel's primitive structural form $1/(1+(\omega\tau)^{2})$ is correctly identified in all four domains, the substrate-ontology evidence accumulates from cross-domain reach. If the kernel's form differs across domains, the substrate-level identification needs revision.

\subsection{Distinctive predictions versus standard proposals}

\begin{center}
\footnotesize
\begin{tabular}{p{0.18\textwidth}p{0.13\textwidth}p{0.13\textwidth}p{0.13\textwidth}p{0.10\textwidth}p{0.18\textwidth}}
\toprule
\textbf{Prediction} & \textbf{Modified dispersion} & \textbf{Lattice cutoff} & \textbf{Naturalness} & \textbf{TCC} & \textbf{ED V5 + DCGT} \\
\midrule
Specific cutoff form & Depends on $f(k\ell_P)$ & Lattice-spacing-dep. & No specific prediction & Cosmo-only & $1/(1+(\omega\tau)^{2})$ FORM-FORCED \\
Cross-platform consistency & Not predicted & Not predicted & Not predicted & Not predicted & $\tau_{V5}$ identification across domains \\
Late-stage PBH evaporation & Mod.-disp.-dep. & Lattice-cutoff-dep. & Robust prediction & No prediction & Planck-mass remnant + spectrum cutoff \\
Connection to soft matter & None & None & None & None & Maxwell relaxation time identification \\
\bottomrule
\end{tabular}
\end{center}

\section{Discussion}

\subsection{The structural-resolution claim}

The trans-Planckian problem of standard semiclassical Hawking has been recognized for over thirty years [1] and addressed through multiple phenomenological approaches. Each approach contributes structural insight; none has produced a derivation-level resolution from a substrate-level ontology.

The Event Density framework's V5 + DCGT account is structurally distinct. The V5 kernel is a substrate primitive existing for reasons independent of the trans-Planckian problem (cross-domain duties in soft matter, BH information, substrate-Unruh). The DCGT bridge is the framework's substrate-to-continuum machinery. Their combination produces a substrate-level UV cutoff at the Planck scale as a \emph{consequence} of substrate ontology, not as a feature added to address the trans-Planckian concern.

This is what we call a \emph{substrate-level resolution}. It differs from phenomenological regularization in that the regulator is not introduced for the problem; it emerges from the substrate-level account that produces the framework's other structural commitments.

\subsection{What this does not claim}

The substrate-level resolution does \emph{not} claim:
\begin{itemize}[noitemsep]
\item That standard semiclassical Hawking is wrong. It correctly describes the empirically validated content (thermal spectrum at $T_H$ at observable scales).
\item That ED is the unique substrate-level framework capable of producing such a resolution. Other substrate ontologies might in principle produce analogous resolutions.
\item That all features of standard Hawking are reproduced exactly. First-subleading-order corrections at $(\ell_P/M)^{2}$ distinguish ED from strict semiclassical Hawking, particularly at extreme scales.
\end{itemize}

What it \emph{does} claim is structural: the substrate ontology produces the trans-Planckian regulator as a derived consequence rather than as a phenomenological addition.

\subsection{Implications for substrate-ontology evidence}

The cross-platform reach of the V5 kernel --- operating in soft-matter rheology, Hawking radiation, BH information, and substrate-Unruh contexts [16] --- provides structural evidence for the substrate ontology. A phenomenological framework would have separate cutoff parameters for each context; ED has one substrate primitive (V5) doing all jobs, with the cutoff scale $\tau_{V5}$ inheriting from the physical context.

The trans-Planckian resolution illustrates this pattern. The same V5 kernel that produces Maxwell viscoelastic memory in soft matter at nanosecond timescales also produces the trans-Planckian cutoff at the Planck scale. The cross-domain reach is structural evidence that V5 is a primitive substrate object, not a domain-specific phenomenological regulator.

\subsection{Open questions}

The substrate-level resolution leaves several questions for future work:

\begin{enumerate}[noitemsep]
\item \textbf{Closed-form derivation of $\tau_{V5}$ at the gravitational scale.} The framework identifies $\tau_{V5} = \ell_P/c$ via dimensional analysis; closed-form derivation from substrate-microscopic V1 + V5 details would tighten the identification.
\item \textbf{The V5 kernel shape function $\psi$.} Different shape functions produce different higher-order corrections; precision tests could distinguish.
\item \textbf{Analog Hawking precision tests.} Current analog experiments confirm spectral form at moderate frequencies; precision measurements at the cutoff scale would test ED's specific functional form against modified-dispersion alternatives.
\item \textbf{Connection to the Standard Model UV completion question.} The V5 substrate cutoff regulates trans-Planckian Hawking modes; whether it also regulates the Standard Model's UV-completion concerns at the Planck scale is open.
\end{enumerate}

These open items do not undermine the substrate-level resolution but identify directions for continued development.

\section{Conclusions}

The trans-Planckian problem of standard semiclassical Hawking radiation arises because the standard derivation traces observed-at-infinity modes back to arbitrarily-blueshifted near-horizon proper frequencies. The standard derivation thus implicitly requires standard QFT to apply at and beyond the Planck scale, where structural concerns about quantum-gravity modification arise.

The Event Density framework provides a substrate-level resolution. The V5 finite-memory cross-chain kernel --- a primitive substrate object with characteristic memory time $\tau_{V5} = \ell_P/c$ at the gravitational scale (FORCED via T19 + substrate dimensional analysis) --- does not coherently mediate cross-chain correlations at frequencies $\omega \gg c/\ell_P$. Combined with the Diffusion Coarse-Graining Theorem's hydrodynamic-window closure when proper wavelengths approach $\ell_P$, the substrate provides a natural UV cutoff at the Planck frequency. The substrate cannot instantiate arbitrarily-blueshifted modes near the horizon; the trans-Planckian divergence of standard semiclassical Hawking is a continuum artifact that vanishes at the substrate level.

The V5-modulated Hawking spectrum is:
\[
N_{ED}(\omega) = \frac{1}{e^{\omega/T_H} - 1} \cdot \frac{1}{1 + (\omega\tau_{V5})^{2}}
\]
finite at all frequencies, reducing to standard Hawking at observable scales, and cut off at the Planck frequency.

The framework's resolution is structurally distinct from standard-physics approaches:
\begin{itemize}[noitemsep]
\item Modified-dispersion approaches postulate UV modification of the dispersion relation.
\item Lattice-cutoff approaches postulate spacetime discreteness at the Planck scale.
\item Naturalness arguments invoke empirical robustness without supplying a mechanism.
\item TCC addresses cosmological inflation, not the BH trans-Planckian problem.
\end{itemize}

ED's V5 + DCGT mechanism produces the cutoff as a \emph{derived consequence} of substrate ontology rather than as a phenomenological addition. The V5 kernel exists for reasons independent of the trans-Planckian problem (cross-domain duties in soft matter, BH information, and substrate-Unruh). The trans-Planckian regulation is a cross-platform application of the same substrate primitive.

Falsifiable predictions distinguishing the framework include the V5-derived cutoff form $1/(1+(\omega\tau)^{2})$ in analog Hawking experiments, the Planck-frequency cutoff in primordial-BH late-stage evaporation, possible high-energy photon dispersion at order $(\omega/\omega_c)^{2}$, and cross-platform identification of $\tau_{V5}$ across the framework's V5-mediated phenomena.

The substrate-level resolution illustrates a general methodological pattern: substrate primitives composed with domain-specific substrate context produce continuum-level observables in apparently unrelated domains, providing structural evidence for the substrate ontology through cross-domain reach.

\section*{References}

\begin{enumerate}[noitemsep,leftmargin=*]
\item Jacobson, T. ``Black-Hole Evaporation and Ultrashort Distances.'' \textit{Physical Review D} \textbf{44}, 1731--1739 (1991).
\item Unruh, W. G. ``Sonic Analog of Black Holes and the Effects of High Frequencies on Black Hole Evaporation.'' \textit{Physical Review D} \textbf{51}, 2827--2838 (1995).
\item 't Hooft, G. ``Dimensional Reduction in Quantum Gravity.'' In: \textit{Salamfest 1993}, 0284--296 [arXiv:gr-qc/9310026] (1993).
\item Helfer, A. D. ``Do Black Holes Radiate?'' \textit{Reports on Progress in Physics} \textbf{66}, 943--1008 (2003).
\item Bedroya, A., Vafa, C. ``Trans-Planckian Censorship and the Swampland.'' \textit{Journal of High Energy Physics} \textbf{2020} (9), 123 (2020).
\item 't Hooft, G. ``On the Quantum Structure of a Black Hole.'' \textit{Nuclear Physics B} \textbf{256}, 727--745 (1985).
\item Hawking, S. W. ``Particle Creation by Black Holes.'' \textit{Communications in Mathematical Physics} \textbf{43}, 199--220 (1975).
\item Hawking, S. W. ``Black Hole Explosions?'' \textit{Nature} \textbf{248}, 30--31 (1974).
\item Steinhauer, J. ``Observation of Quantum Hawking Radiation and Its Entanglement in an Analogue Black Hole.'' \textit{Nature Physics} \textbf{12}, 959--965 (2016).
\item Drori, J., Rosenberg, Y., Bermudez, D., Silberberg, Y., Leonhardt, U. ``Observation of Stimulated Hawking Radiation in an Optical Analogue.'' \textit{Physical Review Letters} \textbf{122}, 010404 (2019).
\item Proxmire, A. \textit{Event Density: One Substrate, Three Domains.} April 2026.
\item Proxmire, A. \textit{Event Density Foundations: A Unified Substrate Architecture for Quantum, Fluid, Gauge, and Gravitational Dynamics.} April 2026.
\item Proxmire, A. \textit{Theorem 19: Newton's Law from Substrate Holographic Counting and the Identification of $\ell_P$.} April 2026.
\item Proxmire, A. \textit{Theorem 18: V1 Kernel Retardation and the Kernel-Level Arrow of Time.} April 2026.
\item Proxmire, A. \textit{The Diffusion Coarse-Graining Theorem: Substrate-to-Continuum Bridge for Canonical-ED Dynamical Content.} April 2026.
\item Proxmire, A. \textit{The V5 Kernel as a Cross-Domain Substrate Mechanism.} May 2026.
\item Proxmire, A. \textit{Arc Hawking H-5: Information Correlations and Substrate-Level Page Curve Structure.} May 2026.
\item Proxmire, A. \textit{Walkthrough: From Primitives to the Substrate-Unruh Effect.} May 2026.
\item Proxmire, A. \textit{Arc Hawking H-1: Spectral Form and Temperature from V5 Cross-Chain Correlations.} May 2026.
\item Proxmire, A. \textit{Walkthrough: From Primitives to Hawking Radiation.} May 2026.
\item Proxmire, A. \textit{Arc Hawking H-8: Higher-Order Resummation and the Late-Time Evaporation Endpoint.} May 2026.
\item Almheiri, A., Marolf, D., Polchinski, J., Sully, J. ``Black Holes: Complementarity or Firewalls?'' \textit{Journal of High Energy Physics} \textbf{2013} (2), 062 (2013).
\item Maldacena, J., Susskind, L. ``Cool Horizons for Entangled Black Holes.'' \textit{Fortschritte der Physik} \textbf{61}, 781--811 (2013).
\item Hawking, S. W., Perry, M. J., Strominger, A. ``Soft Hair on Black Holes.'' \textit{Physical Review Letters} \textbf{116}, 231301 (2016).
\item Almheiri, A., Engelhardt, N., Marolf, D., Maxfield, H. ``The Entropy of Bulk Quantum Fields and the Entanglement Wedge of an Evaporating Black Hole.'' \textit{JHEP} \textbf{2019} (12), 063 (2019).
\item Mattingly, D. ``Modern Tests of Lorentz Invariance.'' \textit{Living Reviews in Relativity} \textbf{8}, 5 (2005).
\item Brout, R., Massar, S., Parentani, R., Spindel, P. ``Hawking Radiation Without Trans-Planckian Frequencies.'' \textit{Physical Review D} \textbf{52}, 4559--4568 (1995).
\item Polchinski, J. ``String Theory and Black Hole Complementarity.'' In: \textit{String Theory} (1999).
\item Brout, R., Massar, S., Parentani, R., Spindel, P. ``A Primer for Black Hole Quantum Physics.'' \textit{Physics Reports} \textbf{260}, 329--446 (1995).
\item Proxmire, A. \textit{Substrate-Level Resolution of the Black Hole Information Paradox.} May 2026.
\item The full Event Density corpus is available at \url{https://github.com/allen-proxmire/event-density}.
\end{enumerate}

\end{document}
